\begin{thebibliography}{99}

\bibitem{anr2023mobilex}
AGENCE NATIONALE DE LA RECHERCHE (ANR). Règlement général du Challenge MOBILEX [en ligne]. Paris : ANR, 2023. Disponible sur : <\href{https://anr.fr/fileadmin/aap/2023/aap-mobilex-2023-reglement-general.pdf}{https://anr.fr/fileadmin/aap/2023/aap-mobilex-2023-reglement-general.pdf}> (consulté le 16.03.2026).

\bibitem{thrun2009robotic}
THRUN S. Robotic Mapping: A Survey. In: Exploring Artificial Intelligence in the New Millenium. San Francisco: Morgan Kaufmann, 2009.

\bibitem{lombard2020stochastic}
LOMBARD C. D., VAN DAALEN C. E. Stochastic triangular mesh mapping: A terrain mapping technique for autonomous mobile robots. Robotics and Autonomous Systems, 127, 103449, 2020.

\bibitem{jaspers2017multi}
JASPERS H., HIMMELSBACH M., WUENSCHE H. J. Multi-modal local terrain maps from vision and lidar. In: 2017 IEEE Intelligent Vehicles Symposium (IV), pp. 1119-1125. IEEE, 2017.

\bibitem{guerrero2015lidar}
GUERRERO J. A., JAUD M., LENAIN R., ROUVEURE R., FAURE P. Towards LIDAR-RADAR based terrain mapping. In: 2015 IEEE International Workshop on Advanced Robotics and its Social Impacts (ARSO), pp. 1-6. IEEE, 2015.

\bibitem{yoloworld2024}
TINGCHENG T., XIZHOU Z., LEHUI L., XIAOHUI X., JIANWEI Y., LEWEN W., XIAOHUA Z., JIANFENG G., BIN X., YUJUN L., XIAOGANG W. YOLO-World: Real-Time Open-Vocabulary Object Detection. arXiv preprint arXiv:2402.09287, 2024.

\bibitem{efficientsam2024}
ZHU X., YANG F., KIRILLOV A., GIRSHICK R., DOLLÁR P., HE K. EfficientSAM: Leveraged Masked Image Pretraining for Efficient Segment Anything. arXiv preprint arXiv:2401.06397, 2024.

\bibitem{segmentanything2023}
KIRILLOV A., MINTUN E., RAVI N., MAITIN-SHEPARD J., GORDON G., SHARMA P., RYALI C., DASARI N., LIAO H., CHEN X., ZHU X., GABEUR V., LIU P., YANG F., DAMON L., DOLLÁR P., GIRSHICK R. Segment Anything. arXiv preprint arXiv:2304.02643, 2023.

\bibitem{hart1968formal}
HART P. E., NILSSON N. J., RAPHAEL B. A formal basis for the heuristic determination of minimum cost paths. IEEE Transactions on Systems Science and Cybernetics, 1968, vol. 4, no. 2, pp. 100-107.

\bibitem{dijkstra1959note}
DIJKSTRA E. W. A note on two problems in connexion with graphs. Numerische Mathematik, 1959, vol. 1, pp. 269-271.

\bibitem{prm1996}
KAVOUSSIANOS J. C., LATOMBE J.-C. Probabilistic roadmaps for path planning in high-dimensional configuration spaces. IEEE Transactions on Robotics and Automation, 1996, vol. 12, no. 4, pp. 566-580.

\bibitem{rrt2000}
LAVALLE S. M., KUFFNER J. J. Rapidly-exploring random trees: A new tool for path planning. In: Proceedings of the IEEE International Conference on Robotics and Automation, 2000, pp. 1297-1302.

\bibitem{rrtstar2011}
KARAMOUZAS I., FRAZOLLI E. RRT*: Optimal Rapidly-exploring Random Trees. In: Proceedings of the International Workshop on the Algorithmic Foundations of Robotics, 2011, pp. 1-10.

\bibitem{deremetz}
DEREMETZ M., LENAIN R., COUVENT A., THUILOT B., CARIOU C. A generic control framework for mobile robots edge following. International Conference Informatics in Control, Automation and Robotics, à paraître.

\bibitem{fox1997dynamic}
FOX D., BURGARD W., THRUN S. The dynamic window approach to collision avoidance. IEEE Robotics \& Automation Magazine, 4(1), pp. 23-33, 1997.

\bibitem{akmandor20203d}
AKMANDOR N. Ü., PADIR T. A 3D reactive navigation algorithm for mobile robots by using tentacle-based sampling. In: 2020 Fourth IEEE International Conference on Robotic Computing (IRC), pp. 9-16. IEEE, 2020.

\bibitem{ozdemir2018follow}
ÖZDEMIR A., SEZER V. Follow the gap with dynamic window approach. International Journal of Semantic Computing, 12(01), pp. 43-57, 2018.

\bibitem{rummelhard2014probabilistic}
RUMMELHARD L., NÈGRE A., PERROLLAZ M., LAUGIER C. Probabilistic Grid-based Collision Risk Prediction for Driving Application. Experimental Robotics, pp. 821–834, 2014.

\bibitem{laconte2021lambda}
LACONTE J. Lambda-Field: a Novel Framework For Risk Assessment In Occupancy Grids. Thèse Robotique. Clermont-Ferrand: Université Clermont Auvergne, 2021.

\bibitem{laconte2021novel}
LACONTE J., KASMI A., POMERLEAU F., CHAPUIS R., MALATERRE L., et al. A Novel Occupancy Mapping Framework for Risk-Aware Path Planning in Unstructured Environments. Sensors, 21(22), pp. 7562, 2021.

\bibitem{benrabah2023dual}
BEN RABAH M., RANDRIAMIARINTSOA E., MORCEAUX J., CHAPUIS R., AUFRERE R. Dual occupancy and knowledge maps management for optimal traversability risk analysis. IEEE Fusion, 2023.

\bibitem{hill2022online}
HILL A., et al. Online Tuning of Control Parameters for Off-Road Mobile Robots: Novel Deterministic and Neural Network-Based Approaches. IEEE Robotics \& Automation Magazine, 30(3), pp. 44-55, 2022.

\bibitem{RST}
LENGAGNE S. Correction des systèmes linéaires échantillonnés [En ligne]. Support de cours. Polytech Clermont. [s.d.]. Disponible sur : <\href{https://moodle2024.uca.fr/pluginfile.php/691566/mod_resource/content/2/Auro3_Seance.pdf}{https://moodle2024.uca.fr/pluginfile.php/691566/mod\_resource/content/2/Auro3\_Seance.pdf}> (Consulté le 20.05.2026).

\bibitem{purepursuit}
THOMASFERMI. Algorithms for Automated Driving: Pure Pursuit. [En ligne]. Disponible sur : <\href{https://thomasfermi.github.io/Algorithms-for-Automated-Driving/Control/PurePursuit.html}{https://thomasfermi.github.io/Algorithms-for-Automated-Driving/Control/PurePursuit.html}> (Consulté le 17.06.2026).

\bibitem{LOAM}
J. Zhang and S. Singh, “Loam : Lidar odometry and mapping in real-time,” in Proceedings of Robotics. Disponible sur : <\href{https://www.researchgate.net/publication/282704722_LOAM_Lidar_Odometry_and_Mapping_in_Real-time}{https://www.researchgate.net/publication/282704722_LOAM_Lidar_Odometry_and_Mapping_in_Real-time}> (Consulté le 28.07.2026).

\bibitem{kitware_slam}
Kitware, “Lidar SLAM : Spotlight on Kitware’s open-source library.” Disponible sur : <\href{https://www.kitware.com/lidar-slam-spotlight-on-kitwares-open-source-library/}{https://www.kitware.com/lidar-slam-spotlight-on-kitwares-open-source-library/}>, 2022. Blog officiel.

\bibitem{dlio_slam}
K. Chen, R. Nemiroff, and B. T. Lopez, “Direct lidar-inertial odometry : Lightweight lio with continuous-time motion correction,” in IEEE International Conference on Robotics and Automation (ICRA), 2023. Disponible sur : <\href{https://arxiv.org/pdf/2203.03749}{https://arxiv.org/pdf/2203.03749}>.

\bibitem{coarse_to_fine}
\href{https://www.researchgate.net/publication/336438983_An_Offline_Coarse-To-Fine_Precision_Optimization_Algorithm_for_3D_Laser_SLAM_Point_Cloud}{An Offline Coarse-To-Fine Precision Optimization Algorithm for 3D Laser SLAM Point Cloud}

\bibitem{gicp}
A. Segal, D. Haehnel, and S. Thrun, “Generalized-icp,” in Proceedings of Robotics : Science and Systems (RSS), 2009.

\bibitem{Kabsch}
KABSCH W. A solution for the best rotation to relate two sets of vectors. Disponible sur : <\href{https://en.wikipedia.org/wiki/Kabsch_algorithm}{https://en.wikipedia.org/wiki/Kabsch_algorithm}> (Consulté le 28.06.2026).

\bibitem{eband}
QUINLAN S., Khatib O. Elastic Bands: Connecting Path Planning and Control. In: Proceedings of the IEEE International Conference on Robotics and Automation (ICRA), 1993. Disponible sur : <\href{https://khatib.stanford.edu/publications/pdfs/Quinlan_1993_ICRA.pdf}{https://khatib.stanford.edu/publications/pdfs/Quinlan_1993_ICRA.pdf}>.

\end{thebibliography}  